\def\velkostpisma{10}
\def\velkostriadku{12.5}
\input makra.tex % nacitanie Ivanom pripravenych nastaveni a prikazov
\hyphenation{star-šov-stvo} % rozdelenie slov na konci riadku, treba tu uviest slova, ktore sam nepozna

\spravodaj{6}{2026}


\clanok{Verš na mesiac}
V~júni sa budeme učiť nový verš, ktorý dostal pre rok~2026 náš spevokol:

„Nech ťa neopustia milosrdenstvo a pravda; priviaž si ich na svoje hrdlo; napíš si ich na tabuľu svojho srdca, a tak nájdi milosť a výborný rozum v~očiach Boha i ľudí.“ (Pr 3,3-4)


\clanok{{\it ČO NÁS ČAKÁ V~JÚNI}}


\clanok{Dejiny židovstva a židovských komunít na Slovensku 5. -- 6.~6.}
Podujatie bolo z bezpečnostných dôvodov ZRUŠENÉ.


\clanok{Rozhovory na Palisádach}
V~sobotu dňa 6.~6. o~9.30~hod. na Palisádach prebehne ďalší z~„Rozhovorov na Palisádach“. Rečníkom bude Pavel Makovíni s~témou „Biblia a peniaze“.
Pozvánka s~podrobnejšími informáciami je aj vo vašich e-mailových schránkach a na nástenke. Ak vás táto téma zaujíma, ste srdečne vítaní; pozvite aj svojich priateľov.


\clanok{Spoločný obed v~SED}
V~júni máme naplánovaný ďalší náš spoločný obed. Stretneme sa v~nedeľu 14.~6. o~12.00 v~priestoroch SED na Partizánskej. Prihlasovať sa nemusíte. Prineste si svoj obed a strávte tento čas pri spoločnom stolovaní a rozhovoroch. Nezabudnite drobnú hotovosť do pokladničky na kávu.


\clanok{Filipíny, krajina tisícov ostrovov, viery a nádeje}
Prednáška s~diskusiou o~misii očami cestovateľky a misionárky našej s.~Rút Krajčiovej, ktorá bola súčasťou medzinárodného tímu. Tím vykonával misiu na ulici s~večerným programom a rozdávaním potravinových balíčkov pre chudobných.
Viac sa dozviete v~nedeľu 14.~6. o~18.00~hod. v~našej modlitebni na Palisádach. Všetci ste srdečne vítaní či už s~priateľmi, susedmi, známymi.


\clanok{Oblastné stretnutie sestier bratislavskej oblasti v~Bernolákove}
V~sobotu 20.~6. sa uskutoční oblastné stretnutie sestier, ktoré sa bude konať v~modlitebni BJB~Bernolákovo. Sestry budú mať svoj program od~9.00 do~16.00~hod. Témou „Víťazný život v~Kristovi“ ich bude sprevádzať s.~Danka Hanesová.


\clanok{{\it PRÁZDNINY}}


\clanok{Zborové tábory}

\cast{Denný detský tábor pre deti od~6 do~10 rokov, 6. -- 10.~7}
Tábor sa bude konať na Súľovskej. Na príprave sa podieľajú bratia kazatelia a pracovníci s~deťmi z~nášho zboru a z~Connectu. Pripravujeme pestrý program a túžime, aby bol pre deti zábavný a osviežujúci, ale zároveň poučný a inšpirujúci.
Prihlasovanie je už uzavreté, prípadné dopyty môžete adresovať s.~Mirke Kešjarovej na \email{mkesjarova@gmail.com}.

\cast{Mládežnícko-dorastový tábor, 19.~--~25.~7.}
Mládež a dorast bude mať svoj tábor v~Novej Lehote. Podrobnejšie informácie dostali táborníci a rodičia e-mailom. Prihlasovanie bolo uzavreté 31.~5.; prípadné otázky vám zodpovie Angie na \email{angela.vrablova@gmail.com}.

\cast{Zborový tábor, 9.~--~15.~8.}
Tábor sa bude konať v~stredisku DM Prameň, Píla Častá. Počas neho bude opäť možnosť verejne vyjadriť svoju vieru v~Pána Ježiša v~krste, ktorý sa uskutoční vo štvrtok 13.~8.
Záujemcovia, prosím, prihlasujte sa u~br.~J.~Szőllősa.
Viac informácií o~tábore, ako aj prihlasovanie (do 30.~6., aj so zálohou), získate kliknutím na link v~e-maile alebo osobne a písomne u~s.~Diany Dzuriakovej \email{diana@dzuriak.com}.


\clanok{Hry namiesto striel -- tábor pre 400~detí v~Černihive}
Naše deti sa z~Božej milosti môžu každoročne zúčastňovať, relaxovať a bezpečne učiť v~rôznych táboroch. Na Ukrajine je to posledné štyri roky zložitejšie... Napriek vojne organizujú tábory pre deti. Tento rok sa z~nášho zboru BJB Palisády do neho zapoja M.~Simon a I.~Kohút.
Modlime sa za nich, za celý organizačný tím, priebeh tábora a najmä za deti v~ňom. Dozvedieť sa viac a finančne podporiť deti na tábore môžete na tejto adrese: \ulink[https://donio.sk/hry-namiesto-striel]{donio.sk/hry-namiesto-striel}. Ďakujeme.


\clanok{{\it OHLIADNUTIE SA ZA MÁJOM}}


\clanok{Správy zo staršovstva}
Staršovstvo sa v~mesiaci máj stretlo dvakrát v~duchu modlitby a spoločného premýšľania nad životom zboru a riešenia tém, ktoré sú dôležité pre jeho ďalšie smerovanie. Venovali sme sa informáciám z~KDZ, príprave zborového členského zhromaždenia, voľbám do orgánov BJB, rozhlasovému nahrávaniu, Chvojnici a ďalším pripravovaným aktivitám zboru.

Jednou z~tém bola budúcnosť služby brata Filipa Barkócziho v~našom zbore. Staršovstvo hovorilo o~možných formách ďalšej spolupráce, finančných súvislostiach a o~neustávajúcej potrebe otvorenej diskusie so zborom.

Výraznou témou bola aj spolupráca s~ukrajinským zborom. Staršovstvo vyhodnotilo spoločné bohoslužby veľmi pozitívne a hovorilo o~možnostiach a túžbe ďalšieho rozvíjania vzťahov. Zároveň sa venovalo aj praktickým otázkam spoločného používania priestorov.

Hovorilo sa tiež o~tom, ako môže byť Súľovská lepšie využitá pre zborové aktivity, mládež, dorast, rodiny či víkendové stretnutia.

Zároveň sme sa venovali aj obnove fasády modlitebne, webu a sociálnym sieťam zboru, novému logu, pripravovanému táboru v~Černihive, prezentácii misijnej cesty sestry Rút Krajčiovej na Filipíny a ďalším praktickým otázkam života zboru.

Prosíme vás o~modlitby za múdrosť pri rozhodovaní, za pokojné a pravdivé rozhovory, jednotu zboru a za všetky služby, ktoré máme ako zbor pred sebou.
\autor{Filip Barkóczi}


\clanok{Večer chvál}
V~piatok 8.~5. sa uskutočnil Večer chvál, ktorý viedol slovom Dávid Chuchút a piesňami Jožko Smutný. Spevom okrem neho slúžili Lenka Pribulová, Adam Alexaj aj s~gitarou, na husliach David Pototskyi a na klavíri doprevádzala Tamarka Smutná. Počas večera zazneli okrem piesní aj Dávidove slová o~vyslobodení, ktorými nadviazal na sviatok 8.~5.

Večer chvál naplnil naše srdcia vďakou a chválou Pána Ježiša, za to aký je a čo pre nás urobil. Túto vzácnu hodinku sme mohli stráviť s~ľudmi okolo nás chváliac Boha a uvedomujúc si Jeho milosť.
\autor{Tamarka Smutná}


\clanok{Deň otvorených dverí v~BJB Palisády}
V~sobotu 7.~5. sa na Súľovskej uskutočnil Deň otvorených dverí, ktorý sme pripravili v~spolupráci s~Connectom. Bolo to požehnané popoludnie, počas ktorého sme mohli stráviť spoločný a neformálny čas.
Deti sa mohli zapojiť do rôznych hier, užiť si skákací hrad a spestrením bol aj kúzelník. Popri programe vznikal priestor na rozhovory, spoznávanie sa a budovanie vzťahov.
Sme vďační Pánu Bohu za príjemnú atmosféru, za všetkých, ktorí prišli, aj za tých, ktorí sa podieľali na príprave a priebehu celého popoludnia.
\autor{Filip Barkóczi}


\clanok{Stretnutie sestier}
Dostala som knižku „Priateľstvo žien“. Premýšľam, čo nové by som sa ja mohla dozvedieť o~priateľstve žien? Príjemný pohľad na rozdiel medzi vzťahmi, ktoré pestujú muži a ktoré ženy. Občas príjemne úsmevné, občas inšpirujúce a naplnené múdrosťou.

Krátko na to máme stretnutie sestier. Skupinka žien, staršie, mladšie i tie najmladšie, jedno malé dievčatko. U~Beátky to má vždy iné čaro. Krásne, oku lahodiace priestory, kde sa s~viditeľnou láskou starajú o~strieborné hlávky.
Sedíme v~tých pôvabných priestoroch na pohodlných kreslách a môj zrak prechádza po jednotlivých ženách, ich príbehoch, radostiach a starostiach. Skrytých aj viditeľných príbehoch. Vidím dlhoročné priateľstvá aj rodinné prepojenia, nové kamarátstva a vidím Toho, ktorý nás všetky spája.
Vtedy začína Beatkin príbeh o~živej vode pre naše životy. Živá voda ako modlitebné stíšenie, ako chvíle pri Božom slove, ako pieseň, ktorá k~nám hovorí, ako človek, cez ktorého prichádza láska. Živá voda, ktorá nám dáva hodnotu a pokoj v~Ježišovi, ktorá nám dáva pokoj v~srdci aj silu odpúšťať.
Pripomíname si popraskané cisterny, ktoré síce dávajú vodu, tá však nie je živá -- poskytuje len krátkodobé utíšenie smädu vo vlastnom výkone, v uznaní od ľudí, v~mnohých vzťahoch, v porovnávaní i v túžbe po dokonalosti a vo vlastnej dokonalosti.

I~takto vyzerá priateľstvo žien: priestor, kde sa môžeme zdieľať a povzbudzovať, pripomínať si dobré veci, učiť sa nové a niesť sa vzájomne na modlitbách.

Drahý Pane, prosím, aby sme vedeli rozoznať a nájsť tú správnu studňu, aby sme mohli ponúknuť živú vodu aj ľuďom okolo seba.
\autor{Slávka Kráľová}


\clanok{Letnice na Chvojnici}
Po roku sme sa opäť vydali na Chvojnicu. Tradične na Letnice ako to bolo po mnohé minulé roky, kedy sme túto slávnosť zoslania Ducha Svätého prežívali s~domácimi Chvojničanmi. Pár rokov sme už sami a v~kostolíku už každú nedeľu neslúžime domácim.
Napriek tomu sa tam stále radi vraciame a vždy sa vyformuje aspoň malá skupinka, ktorá rada oslávi sviatok na Chvojnici. V~sobotu sme sa venovali upratovaniu chalupy a stodoly, koseniu, aby bolo toto miesto pripravené na letné aktivity.
V~nedeľu nám v~modlitebni slúžil br.~kazateľ Janko a chválospevová skupinka. Pri spoločnom obede, guláši, sme sa zišli asi tridsiati. Pozvaní boli aj domáci Peričkovci a Hyžovci, ktorí nám už roky ochotne pomáhajú pri zveľaďovaní chalupy a jej okolia.
Posedenie pri obede, káve a koláčikoch nám priniesli príjemné chvíle pri vzájomných rozhovoroch. Tak nabudúce opäť o~rok \smiley.
\autor{Katarína Kerekréty}


\clanok{Noc otvorených kostolov}
V~piatok 29.~5. v~rámci akcie „Noc kostolov“ prišli do nášho kostola štyria ľudia. Jeden manželský pár cielene a dvoch ďalších sme oslovili vonku pred bránou. Všetci mali záujem vojsť dovnútra a boli prekvapení jednoduchosťou interiéru našej modlitebne.
Gabi, jedna z~návštevníčok, mala veľa otázok: „Kde máte organ?“, „K~čomu sú tie schodíky pod krížom?“ Keď sme odpovedali, že slúžia pre spevákov, tak sa živo zaujímala aj o~náš spevokol.
Mladého muža z~USA zaujal hlavne kríž a bol rád, že sa ho mohol aj dotknúť. Nikdy sa nezaujímal o~kresťanstvo, a tak to bola krásna príležitosť povedať mu o~Ježišovi.

Bola to krásna možnosť byť s~ľuďmi zvonku, ale aj mať spoločenstvo navzájom s~úžasnými sestrami a bratom \smiley\ -- Vladina~L., Janka~L., Zuzka~P., Erik~V., Mirka~H.

Veľmi sa nám to celé páčilo a určite má zmysel ísť do tejto akcie aj o~rok a možno aj skôr, ako je zvykom. Čo tak otvoriť bránu pre návštevníkov už na jeseň? Verím, že Pán Ježiš sa môže dotknúť srdca niektorého z~tých náhodných návštevníkov.
\autor{Mirka Hovorková}


\clanok{Zbierky v~máji}
\table{lr}{
Na misiu           & 600~€ \cr
Na investičný fond &  13~€ \cr
}
\vskip1em

Pripomíname, že zbierky sa u~nás konajú pravidelne takto:
\vskip-1ex\begitems
* každú 2.~nedeľu v~mesiaci je zbierka venovaná misii a
* každú 4.~nedeľu je zbierka na investície.
\enditems

Prispieť je možné formou zbierky počas 4.~nedieľ v~mesiaci alebo bankovým prevodom na zborový účet, prípadne aj osobne u~s.~Ľubky Kohútovej.

Bankové spojenie: SK36 0900 0000 0000 1147 1836, SWIFT: GIBASKBX

\n  2.	6.	Miriam	KEŠJAROVÁ;
\n  7.	6.	Pavel	KOHÚT;
\n  9.	6.	Samuel	PLETT;
\n 10.	6.	Ján	LAURENČÍK;
\n 15.	6.	Ľubica	HOVORKOVÁ;
\n 15.	6.	Trey	ATKINS;
\n 17.	6.	Juraj	KVAČKA;
\n 19.	6.	Anna	ŠANDOROVÁ;
\n 19.	6.	Oľga	VALCHÁŘOVÁ;
\n 22.	6.	Kristína	HORVÁTIKOVÁ;
\n 25.	6.	Peter	ŽEMBERY;
\n 25.	6.	Marica	ŠČEVLÍKOVÁ;
\n 27.	6.	Sylvia	PRIBULOVÁ;
\n 28.	6.	Jana	PERKNOVSKÁ;
\narodeniny


\program{
\p  1 ; po ; 18.00 ; Modlitebná skupinka / skupinka „Rastieme v~Kristovi“ ;.;;
\p  2 ; ut ; 15.00 ; Biblická hodina pre seniorov (P.~Pivka) ;.;;
\p  3 ; st ;.;;.;;
\p  4 ; št ; 18.00 ; Biblická hodina (F.~Barkóczi) ;.;;
\p  5 ; pi ; 17.30 ; Dorast ;.;;
\p  6 ; so ; 18.00 ; Mládež ;.;;
\p  7 ; ne ;  9.30 ; Bohoslužby (J.~Szőllős + VP) ;.;;
\p  8 ; po ; 18.00 ; Modlitebná skupinka ;.;;
\p  9 ; ut ; 15.00 ; Biblická hodina pre seniorov (P.~Pivka) ;.;;
\p 10 ; st ;.;;.;;
\p 11 ; št ; 18.00 ; Biblická hodina (D.~M.~Chuchút) ;.;;
\p 12 ; pi ; 17.30 ; Dorast ;.;;
\p 13 ; so ; 18.00 ; Mládež ;.;;
\p 14 ; ne ;  9.30 ; Bohoslužby (F.~Barkóczi) ; 12.00 ; Spoločný obed v~SED ;
\p 15 ; po ; 18.00 ; Modlitebná skupinka / skupinka „Rastieme v~Kristovi“ ;.;;
\p 16 ; ut ; 15.00 ; Biblická hodina pre seniorov (P.~Pivka) ;.;;
\p 17 ; st ;.;;.;;
\p 18 ; št ; 18.00 ; Biblická hodina (F.~Barkóczi) ;.;;
\p 19 ; pi ; 17.30 ; Dorast ;.;;
\p 20 ; so ; 18.00 ; Mládež ;.;;
\p 21 ; ne ;  9.30 ; Bohoslužby (D.~M.~Chuchút) ;.;;
\p 22 ; po ; 18.00 ; Modlitebná skupinka ;.;;
\p 23 ; ut ; 15.00 ; Biblická hodina pre seniorov (P.~Pivka) ;.;;
\p 24 ; st ; 17.30 ; Stretnutie sestier ;.;;
\p 25 ; št ; 18.00 ; Biblická hodina (D.~M.~Chuchút) ;.;;
\p 26 ; pi ; 17.30 ; Dorast ;.;;
\p 27 ; so ; 18.00 ; Mládež ;.;;
\p 28 ; ne ;  9.30 ; Bohoslužby (Mário Šulka) ;.;;
\p 29 ; po ; 18.00 ; Modlitebná skupinka / skupinka „Rastieme v~Kristovi“ ;.;;
\p 30 ; ut ; 15.00 ; Biblická hodina pre seniorov (P.~Pivka) ;.;;
}

\tiraz
\bye
